% Journal of Intelligent Transportation Systems submission manuscript.
% Compile the named version with: latexmk -pdf manuscript.tex
% Compile the anonymous version with: latexmk -pdf anonymous.tex
\documentclass[largeformat]{interact}

\usepackage[T1]{fontenc}
\usepackage[utf8]{inputenc}
\usepackage{amsmath,amssymb,amsthm}
\usepackage{booktabs,array}
\usepackage{graphicx}
\usepackage[caption=false]{subfig}
\usepackage{adjustbox}
\usepackage[ruled,vlined,linesnumbered]{algorithm2e}
\usepackage{tikz}
\usetikzlibrary{shapes.geometric,arrows.meta,positioning,fit,backgrounds,calc,patterns,decorations.pathreplacing}
\usepackage{pgfplots}
\pgfplotsset{compat=1.18}
\usepgfplotslibrary{fillbetween}
\usepackage[natbibapa]{apacite}
\setlength\bibhang{12pt}
\renewcommand\bibliographytypesize{\fontsize{10}{12}\selectfont}
\usepackage[hidelinks]{hyperref}
\usepackage{microtype}

% The Interact class reserves \doi for article metadata.  Rebind it for
% apacite's bibliography DOI links after hyperref has been loaded.
\renewcommand{\doi}[1]{\href{https://doi.org/#1}{https://doi.org/#1}}

\graphicspath{{figures/}{figures/dreamview/}}
\input{figures/ped_macros}

\definecolor{deepblue}{RGB}{0,114,178}
\definecolor{lightblue}{RGB}{187,220,238}
\definecolor{skyblue}{RGB}{86,180,233}
\definecolor{inputyellow}{RGB}{240,228,66}
\definecolor{outputpink}{RGB}{204,121,167}
\definecolor{improvedpink}{RGB}{213,94,0}
\definecolor{lightpink}{RGB}{230,159,0}
\definecolor{deeppink}{RGB}{213,94,0}
\definecolor{bglight}{RGB}{244,244,244}
\definecolor{lightgreen}{RGB}{0,158,115}
\definecolor{lightorange}{RGB}{230,159,0}
\definecolor{dangerred}{RGB}{213,94,0}
\definecolor{vibrantgreen}{RGB}{0,158,115}
\definecolor{vibrantorange}{RGB}{230,159,0}
\definecolor{vibrantred}{RGB}{213,94,0}
\tikzset{
  mod/.style={rectangle,rounded corners=2pt,draw=deepblue,fill=lightblue,minimum height=0.9cm,align=center,font=\small},
  modimp/.style={rectangle,rounded corners=2pt,draw=deeppink,fill=improvedpink!25,minimum height=0.9cm,align=center,font=\small},
  modin/.style={rectangle,draw=deepblue,fill=inputyellow!45,minimum height=0.9cm,align=center,font=\small},
  modout/.style={rectangle,draw=deeppink,fill=outputpink!35,minimum height=0.9cm,align=center,font=\small},
  axx/.style={->,thick,black!80},
  flow/.style={-Stealth,thick,deepblue},
  flowimp/.style={-Stealth,thick,deeppink},
  obsstatic/.style={fill=dangerred!70,draw=dangerred,thick},
  obsdyn/.style={fill=lightorange,draw=dangerred,thick},
  obsbuf/.style={draw=dangerred!70,dashed,thick},
  bandok/.style={fill=lightgreen,opacity=0.35},
  bandbad/.style={fill=dangerred!25,opacity=0.55},
  egopt/.style={fill=deepblue,draw=deepblue},
  reflinestyle/.style={deepblue,very thick}
}

\theoremstyle{plain}
\newtheorem{theorem}{Theorem}[section]
\newtheorem{lemma}[theorem]{Lemma}
\newtheorem{proposition}[theorem]{Proposition}

\DontPrintSemicolon
\SetAlgoCaptionSeparator{.}
\SetAlgorithmName{Algorithm}{Algorithm}{List of Algorithms}
\providecommand{\pp}{percentage points}
\providecommand{\MIKU}{MIKU}
\providecommand{\code}[1]{\texttt{#1}}

\articletype{Research Article}
\title{Interaction-Aware Constraint Interface for Path--Velocity Planning in Dynamic Multi-Obstacle Traffic}

\ifdefined\anonymousversion
  \author{\name{Anonymous authors}\affil{Affiliations withheld for anonymous review}}
  \hypersetup{pdftitle={Interaction-Aware Constraint Interface for Path-Velocity Planning in Dynamic Multi-Obstacle Traffic},pdfauthor={Anonymous}}
\else
  % Identified author metadata.  Keep this file out of the anonymous source set.
\author{%
  \name{JiaWang Liao\textsuperscript{a}, YuFei Hu\textsuperscript{a}, ChaoYang Shi\textsuperscript{a} and ChengJiao Sun\textsuperscript{a}\thanks{CONTACT ChengJiao Sun. Email: jiao1952@126.com}}
  \affil{\textsuperscript{a}School of Computer Engineering, Hubei University of Arts and Science, 296 Longzhong Road, Xiangyang 441053, Hubei, China}}
\hypersetup{pdftitle={Interaction-Aware Constraint Interface for Path-Velocity Planning in Dynamic Multi-Obstacle Traffic},pdfauthor={JiaWang Liao, YuFei Hu, ChaoYang Shi, ChengJiao Sun}}

\fi

\begin{document}
\maketitle

\begin{abstract}
In two-stage path--velocity planning, time-blind obstacle bounds and local passing-side decisions can rule out feasible passages and cause unnecessary stops in dense traffic. We present MIKU, a constraint-generation method that constructs lateral bands from estimated arrival times, selects a sequence of bands across road stations, and maps safe passing windows into bounds for the existing path and speed quadratic programs. For a fixed-section interval model, the widest lateral band can be identified by sorting and scanning obstacle intervals in $O(k\log k)$ time. In 3{,}500 paired structured-roadway simulations spanning seven traffic scenarios, MIKU increased collision-free goal reaching from 62.69\% for an internal time-blind, fixed-side baseline to 77.46\% (paired difference: 14.77 percentage points; 95\% confidence interval: 13.54--15.97). In a separate set of 700 rolling-replanning simulations, success increased from 62.29\% to 72.29\%. Complete-call planning latency had a median of 9.98\,ms and a 99th percentile of 221.80\,ms in the paired benchmark. These findings show that interaction-aware constraint selection improves path--velocity planning under the tested simulation conditions.
\end{abstract}

\begin{keywords}
autonomous driving; path--velocity decomposition; space--time homotopy; dynamic obstacles; motion planning
\end{keywords}

\section{Introduction}
\label{sec:introduction}

In urban mixed-traffic autonomous driving, planners built on path--velocity decomposition emit unnecessary fallback stops because time-blind path bounds and obstacle-wise side choices discard interaction timing before either solver runs. For autonomous-vehicle system integrators and safety engineers, these artificial deadlocks directly degrade ride quality and route efficiency, creating a measurable gap between theoretical feasibility and operational performance in dense traffic. This paper develops MIKU, a constraint reconstruction framework for two-stage quadratic programming path--velocity decomposition architectures, evaluated on 4,000 deterministic interval instances and 3,500 paired simulated scenarios via collision-free goal reaching, progress ratio, jerk RMS, and planning-time quantiles, with degraded-stop rate directly reflecting the operational trade-offs among safety, progress, ride comfort, and real-time computation \citep{gonzalez2016,paden2016,schwarting2018planning}.

Specifically, path--velocity decomposition separates lateral path generation in Frenet station--lateral coordinates $s$--$l$ from longitudinal speed optimization in station--time coordinates $s$--$t$ \citep{kant1986pvd,werling2010,zhou2021dliaps}. Industrial implementations typically solve two structured convex quadratic programs that expose environmental constraints through spatial path boundaries and temporal ST boundaries \citep{fan2018baidu,zhang2020pathqp,wang2022_jfr}. In dense dynamic scenes this separation discards two kinds of information: a path boundary from a current snapshot or a horizon union cannot distinguish an area occupied now from one free at the ego's arrival, and irreversible obstacle-wise side choices can raise the lower and lower the upper path bound in mutually incompatible ways even when a jointly selected passage remains open.

While joint space--time search, mixed-integer optimization, and spatio-temporal corridor methods coordinate spatial and temporal decisions simultaneously \citep{mcnaughton2011lattice,kessler2023,yoon2024stcorridor}, they typically expand state spaces or alter optimization structures, making them incompatible with production two-QP architectures. Consequently, current industrial PVD implementations suffer from three fundamental limitations:
first, time-blind path boundaries merge occupancies across the prediction horizon, losing arrival-conditioned clearance before lateral optimization;
second, greedy, obstacle-wise side choices cannot capture group-level continuous passages, causing lower and upper path bounds to cross and falsely signal a blockage even when viable corridors exist;
and third, existing cross-stage coordination mechanisms cannot transfer spatio-temporal homotopy decisions without modifying downstream convex solvers.
As illustrated in Figure~\ref{fig:teaser_en}, MIKU resolves these challenges directly at the constraint interface by constructing time-dependent lateral bands, selecting spatial and temporal homotopy classes, and mapping them into standard box constraints accepted by the two QPs.

\begin{figure}
\centering
\adjustbox{max width=\textwidth}{\begin{tikzpicture}[x=0.82cm,y=0.82cm,>=Stealth]
\begin{scope}
  \fill[bglight] (-.4,-2.25) rectangle (7.8,2.25);
  \draw[gray,thick] (-.4,2.25)--(7.8,2.25) (-.4,-2.25)--(7.8,-2.25);
  \draw[dashed,gray!60] (-.4,0)--(7.8,0);
  \fill[dangerred!70] (2.7,.45) rectangle (4.3,1.0);
  \fill[dangerred!70] (2.7,-1.0) rectangle (4.3,-.45);
  \draw[dangerred!60,dashed] (2.3,.05) rectangle (4.7,1.4);
  \draw[dangerred!60,dashed] (2.3,-1.4) rectangle (4.7,-.05);
  \node[font=\scriptsize] at (3.5,.73) {$O_a$};
  \node[font=\scriptsize] at (3.5,-.73) {$O_b$};
  \node[deeppink,font=\scriptsize\bfseries] at (5.0,.8) {$R$};
  \node[deepblue,font=\scriptsize\bfseries] at (2.0,-.8) {$L$};
  \draw[deeppink,very thick] (-.4,2.25)--(2.1,2.25)--(2.1,0)--(4.9,0)--(4.9,2.25)--(7.8,2.25);
  \draw[deepblue,very thick] (-.4,-2.25)--(2.1,-2.25)--(2.1,0)--(4.9,0)--(4.9,-2.25)--(7.8,-2.25);
  \node[dangerred,font=\scriptsize\bfseries] at (3.5,1.75) {$l^+\le l^-\;\Rightarrow\;$ blocked};
  \node[font=\small\bfseries] at (3.5,-2.75) {Obstacle-wise greedy choice};
\end{scope}
\begin{scope}[xshift=9cm]
  \fill[bglight] (-.4,-2.25) rectangle (7.8,2.25);
  \draw[gray,thick] (-.4,2.25)--(7.8,2.25) (-.4,-2.25)--(7.8,-2.25);
  \draw[dashed,gray!60] (-.4,0)--(7.8,0);
  \fill[dangerred!70] (2.7,.45) rectangle (4.3,1.0);
  \fill[dangerred!70] (2.7,-1.0) rectangle (4.3,-.45);
  \draw[dangerred!60,dashed] (2.3,.05) rectangle (4.7,1.4);
  \draw[dangerred!60,dashed] (2.3,-1.4) rectangle (4.7,-.05);
  \node[font=\scriptsize] at (3.5,.73) {$O_a$};
  \node[font=\scriptsize] at (3.5,-.73) {$O_b$};
  \node[deepblue,font=\scriptsize\bfseries] at (5.0,.8) {$R$};
  \node[deepblue,font=\scriptsize\bfseries] at (5.0,-.8) {$R$};
  \fill[lightgreen,opacity=.28] (-.4,-2.25) rectangle (7.8,-1.45);
  \draw[deeppink,very thick] (-.4,2.25)--(2.1,2.25)--(2.1,-1.45)--(4.9,-1.45)--(4.9,2.25)--(7.8,2.25);
  \draw[deepblue,dashed,very thick] (-.4,-2.25)--(7.8,-2.25);
  \draw[deepblue,very thick,->] (.3,-1.85)--(2.2,-1.85)--(4.0,-1.80)--(5.8,-1.85)--(7.4,-1.85);
  \node[font=\small\bfseries] at (3.5,-2.75) {Group-level continuous partition};
\end{scope}
\end{tikzpicture}
}
\caption{Greedy obstacle-wise side choices make $l^+(s)$ and $l^-(s)$ cross and trigger a fallback stop (a). MIKU's group-level partition keeps a continuous passage for $\mathcal{T}_{\mathrm{ego}}$ (b).}
\label{fig:teaser_en}
\end{figure}

The principal contributions of this work are as follows:
\begin{enumerate}
  \item \textbf{Arrival-time-dependent spatial projection:} Addressing the first gap of time-blind occupancy merging, we formulate an arrival-time-dependent mapping that projects obstacle occupancies into center-coordinate path bounds based on estimated ego arrival times, ensuring the path planner considers only obstacles present when the vehicle traverses each station. This mechanism is developed in Section~\ref{sec:method_projection_en} and isolated by ablation~A1 in Table~\ref{tab:ablation_results}.
  \item \textbf{Group-level continuous lateral partitioning and fast scan:} Addressing the second gap of greedy obstacle-wise deadlocks, we prove the continuous-partition lemma and maximum-gap theorem in Section~\ref{sec:method_band_en}. These theoretical foundations reduce the exponential $2^k$ side assignments for $k$ obstacles to an exact $O(k\log k)$ scan under the fixed-section interval model, verified against exhaustive search across 4,000 deterministic instances and evaluated by ablations~A3--A4 in Table~\ref{tab:ablation_results}.
  \item \textbf{Layered temporal homotopy graph and QP-preserving transfer:} Addressing the third gap of cross-stage constraint loss, we transfer spatial decisions to longitudinal planning via uncertainty-expanded occupancy tubes, safe passage windows, and a layered temporal homotopy graph in Section~\ref{sec:method_temporal_en}. Timing conditions are formulated as linear box constraints that preserve downstream QP structures and computational efficiency, with empirical verification through ablations~A5--A6 in Table~\ref{tab:ablation_results}.
  \item \textbf{Comprehensive empirical and system-level validation:} We systematically evaluate MIKU across 3,500 paired scenarios with rigorous confidence intervals and ablations, validate closed-loop execution in 700 rolling-replanning episodes, and demonstrate software-interface compatibility with the production Apollo 11.0 autonomous driving system in Sections~\ref{sec:experimental_design} and \ref{sec:results}.
\end{enumerate}

Section~\ref{sec:related_work} positions the interface among joint and decomposed planners; Section~\ref{sec:problem} defines the PVD interface and its failure mechanisms; Section~\ref{sec:method} presents MIKU; Sections~\ref{sec:experimental_design}--\ref{sec:results} cover evaluation and results; Sections~\ref{sec:discussion}--\ref{sec:conclusion} discuss and conclude.

\section{Related work}
\label{sec:related_work}

\subsection{Joint space--time planning}

Joint planners retain geometric and temporal coupling by searching or optimizing an expanded state. Conformal spatiotemporal lattices coordinate position, velocity, and time directly \citep{mcnaughton2011lattice}; instantaneous-analysis and SLT formulations similarly combine lateral and longitudinal decisions in dense traffic \citep{yang2021_ia_planner,hu2023_hybrid_astar_slt,hu2025_dp_nmpc,qiao2025stjoint}. Mixed-integer formulations can encode obstacle-side choices as discrete variables \citep{kessler2023}. These approaches provide strong coordination, but their expanded search spaces, integer decisions, or nonlinear programs differ structurally from the two sequential convex QPs used by many PVD stacks. The finite-grid joint search used later in this paper is therefore a computational reference rather than a claimed continuous optimum.

\subsection{Decomposed and iterative coordination}

PVD planners obtain tractability by optimizing a path first and then a speed profile \citep{kant1986pvd,werling2010}. Convex smoothing and piecewise-jerk optimization make this structure suitable for production-oriented software \citep{fan2018baidu,zhou2021dliaps,zhou2021pathqp}. Temporal optimization, iterative anchoring, candidate pairing, and risk-aware coordination can return information between the two stages \citep{liu2017speed,chen2019em,zhang2023_dp_rcg,han2026_contingency,pathspeedstructured2024}. Their usefulness depends on what information each stage exposes: repeated arrival-time updates cannot recover a lateral alternative that was discarded by an irreversible local side decision.

Safe-corridor methods make the feasible region explicit before continuous optimization. Spatial flight corridors demonstrate how convex sets can bridge discrete search and trajectory optimization \citep{liu2017sfc}, while recent autonomous-driving work extends this idea to spatiotemporal lane-change or Frenet corridors \citep{yoon2024stcorridor,frenetcorridor2025}. MIKU shares the explicit-feasible-region principle but targets a narrower interface question: how to generate mutually consistent SL and ST bounds without replacing the path and speed QPs.

\subsection{Interaction prediction and safety margins}

Prediction-aware planning can account for pedestrian and vehicle interactions that a static obstacle footprint misses \citep{yangbo2023ped,jeong2021ml}. Such predictions remain uncertain, so a planner must distinguish a nominal arrival estimate from the occupancy set used for collision checking. MIKU uses the former to query and rank alternatives, while bounded position and velocity errors expand the latter. Threat-dependent extra margins further allocate scarce lateral space according to collision urgency, overlap, relative motion, object type, and local density; time-to-collision is used only as one normalized urgency factor \citep{minderhoud2001extended}.

Table~\ref{tab:method_classes} summarizes the solver-level distinction. Existing families are effective under their stated assumptions; the unresolved case is a PVD stack that must preserve its downstream QPs while jointly coordinating time-dependent projection, group-level side selection, and speed-stage timing constraints.

\begin{table}
\tbl{Relationship between planning families and the constraint-interface problem.}{%
\begin{tabular}{p{0.19\textwidth}p{0.24\textwidth}p{0.23\textwidth}p{0.24\textwidth}}
\toprule
Family & Space--time coordination & Solver implication & Interface-level gap addressed here \\
\midrule
Time-blind PVD & Sequential path then speed & Two structured QPs & Loses arrival-conditioned occupancy and group choices \\
Iterative PVD & Repeated path--speed updates & Multiple solver calls & Cannot restore discarded lateral classes by timing alone \\
Joint lattice or mixed-integer planning & Coupled state or discrete variables & Expanded search or a different optimizer & Does not preserve the original two-QP interface \\
Safe-corridor planning & Convex feasible sets before optimization & Corridor generator plus continuous solver & Often focuses on a selected corridor rather than coupled spatial and temporal classes \\
MIKU & Spatial and temporal homotopies compiled into bounds & Retains the two downstream QPs & Targets arrival timing, group decisions, and cross-stage constraint transfer \\
\bottomrule
\end{tabular}}
\label{tab:method_classes}
\end{table}

\section{Problem formulation}
\label{sec:problem}

We consider an autonomous vehicle navigating a roadway over planning horizon $T_{\mathrm{pred}}$ in the presence of $k(t)$ static and dynamic obstacles. Given vehicle kinematic states and obstacle motion predictions, our objective is to systematically construct the spatial and temporal constraint interface $(\mathcal{B}_l,\mathcal{B}_s)$ for two-stage path--velocity optimization without altering the structure of downstream convex quadratic programs. Following standard formulations in autonomous vehicle motion planning, we focus on behavioral- and trajectory-level optimization; low-level tracking controllers are assumed to track reference trajectories, while communication delays and low-level powertrain dynamics are decoupled from the upper-level trajectory generation.

\subsection{Path--velocity decomposition interface}

We formulate vehicle motion planning in a road-aligned Frenet coordinate frame defined with respect to a designated smooth reference line. In this frame, the longitudinal coordinate $s$ denotes the arc length along the reference line, representing the vehicle's forward travel distance, while the lateral coordinate $l$ denotes the signed orthogonal displacement from the reference line.

Under the two-stage path--velocity decomposition framework, the complete spatio-temporal motion planning problem is decoupled into two sequential stages:
first, solving for a spatial path profile $l(s)$ that determines lateral maneuvering and obstacle avoidance along the longitudinal station;
second, solving for a longitudinal trajectory profile $s(t)$ that governs the vehicle's forward position, speed, and acceleration over the planning horizon $t \in [0, T_{\mathrm{pred}}]$.

To render these continuous profiles tractable for real-time numerical optimization, the curves are discretized over spatial stations $s_j$ for $j=0,\dots,K$ and temporal steps $t_r$ for $r=0,\dots,N$. Let the discrete path state vector be $\mathbf{x}_p=(l_j,l'_j,l''_j)_{j=0}^{K}$, where $l'_j = \mathrm{d}l/\mathrm{d}s$ and $l''_j = \mathrm{d}^2l/\mathrm{d}s^2$ represent the lateral heading deflection and curvature rate, respectively. Let the discrete longitudinal trajectory state vector be $\mathbf{x}_v=(s_r,v_r,a_r)_{r=0}^{N}$, where $v_r = \dot{s}(t_r)$ and $a_r = \ddot{s}(t_r)$ denote the longitudinal velocity and acceleration. The two planning stages are sequentially solved via convex quadratic programs:
\begin{align}
 \mathbf{x}_p^*={}&\arg\min_{\mathbf{x}_p}
 \frac{1}{2}\mathbf{x}_p^{\mathsf T}H_p\mathbf{x}_p+q_p^{\mathsf T}\mathbf{x}_p,\\[-2pt]
 &\quad\text{s.t. }A_p\mathbf{x}_p=b_p,\quad l_j^-\leq l_j\leq l_j^+, \label{eq:path_qp_en}\\
 \mathbf{x}_v^*={}&\arg\min_{\mathbf{x}_v}
 \frac{1}{2}\mathbf{x}_v^{\mathsf T}H_v\mathbf{x}_v+q_v^{\mathsf T}\mathbf{x}_v,\\[-2pt]
 &\quad\text{s.t. }A_v\mathbf{x}_v=b_v,\quad s_r^{lb}\leq s_r\leq s_r^{ub},\notag\\[-2pt]
 &\qquad 0\leq v_r\leq v_{\max},\quad a_{\min}\leq a_r\leq a_{\max}. \label{eq:speed_qp_en}
\end{align}
Here, $H_p$ and $H_v$ encode trajectory smoothness costs, while $A_p$ and $A_v$ enforce piecewise-jerk kinematic transition equalities between adjacent discretization points. The operational environment is exposed to the solvers strictly through the spatial path boundaries $\mathcal{B}_l=\{[l_j^-,l_j^+]\}_{j=0}^K$ and temporal position boundaries $\mathcal{B}_s=\{[s_r^{lb},s_r^{ub}]\}_{r=0}^N$. MIKU reconstructs these constraint interfaces while preserving the objective matrices, kinematic equalities, and convex QP solver formulations.

\subsection{Center-coordinate obstacle model}

For obstacle $i$, let $O_i(t)=[s_i^-(t),s_i^+(t)]\times[l_i^-(t),l_i^+(t)]$ denote its physical bounding box. With road limits $[l_{\mathrm{road}}^-,l_{\mathrm{road}}^+]$, vehicle width $W_{\mathrm{ego}}$, and roadside margin $d_{\mathrm{road}}$, the feasible interval for the vehicle center is
\begin{equation}
 [L^-,L^+]=[l_{\mathrm{road}}^-+W_{\mathrm{ego}}/2+d_{\mathrm{road}},\ l_{\mathrm{road}}^+-W_{\mathrm{ego}}/2-d_{\mathrm{road}}].
 \label{eq:center_road_en}
\end{equation}
The obstacle occupancy interval is expanded once by half the vehicle width and a threat-dependent buffer $d_{\mathrm{buf},i}$:
\begin{equation}
 [u_i,v_i]=[l_i^- -W_{\mathrm{ego}}/2-d_{\mathrm{buf},i},\ l_i^+ +W_{\mathrm{ego}}/2+d_{\mathrm{buf},i}].
 \label{eq:center_forbidden_en}
\end{equation}
All quantities in Equations~\ref{eq:center_road_en}--\ref{eq:center_forbidden_en} are defined with respect to the vehicle center, ensuring that geometric dimensions are not subtracted repeatedly in downstream modules.

Assigning each of the $k$ active obstacles at a given longitudinal station to the left or right side via decision vector $\mathbf{d}$ yields
\begin{equation}
 l^-(\mathbf{d})=\max\left(L^-,\max_{i\in\mathcal{L}}v_i\right),\qquad
 l^+(\mathbf{d})=\min\left(L^+,\min_{i\in\mathcal{R}}u_i\right),
 \label{eq:center_band_en}
\end{equation}
with empty index sets omitted. The resulting lateral band is feasible when the available clearance satisfies $W(\mathbf{d})=l^+(\mathbf{d})-l^-(\mathbf{d})\geq\varepsilon$ for a minimum width threshold $\varepsilon$.

\subsection{Baseline failure mechanisms and evaluation criteria}

A time-blind path stage uses the full prediction horizon union
\begin{equation}
 \overline{O}_i=\bigcup_{t\in[0,T_{\mathrm{pred}}]}O_i(t),
 \label{eq:prob_union_en}
\end{equation}
treating spatial occupancies across different time instances as simultaneously blocked. Furthermore, greedy obstacle-wise side decisions can push the upper and lower bounds in Equation~\ref{eq:center_band_en} past each other, resulting in $l^+(s)<l^-(s)$ and falsely declaring an artificial blockage even when a collision-free spatio-temporal passage exists.

An episode is deemed successful when the optimized trajectory composed of $\mathbf{x}_p^*$ and $\mathbf{x}_v^*$ remains collision-free throughout $[0,T_{\mathrm{pred}}]$ and reaches $s\geq s_{\mathrm{end}}-1$~m, with the complete operational evaluation metrics detailed in Section~\ref{sec:experimental_design}.

\section{MIKU methodology}
\label{sec:method}

\subsection{Overview}

MIKU reconstructs the two pieces of information lost by time-blind and obstacle-wise PVD. It first estimates the ego arrival time $\tau(s)$ and queries time-dependent obstacle occupancy. It then computes threat-dependent center bounds, groups longitudinally connected obstacles, and selects continuous lateral bands. After the path QP, uncertainty-expanded occupancy tubes define safe passage windows. A layered temporal homotopy graph selects a causally compatible window sequence and writes it into the ST position bounds. Figure~\ref{fig:miku_method_flow_en} summarizes the data flow.

\begin{figure}
\centering
\adjustbox{max width=\textwidth}{\begin{tikzpicture}[x=.92cm,y=.92cm,>=Stealth]
\node[modin,minimum width=2.55cm] (input) at (-0.5,1.05) {Road geometry\\Ego state\\Obstacle predictions};
\node[modimp,minimum width=2.55cm] (occupancy) at (3.35,1.05) {Arrival-time query\\Threat margin\\Robust occupancy};
\node[modimp,minimum width=2.45cm,text width=2.35cm] (spatial) at (6.70,1.05) {Continuous partition\\Top-$K$ spatial classes};
\node[mod,minimum width=1.95cm] (path) at (9.55,1.05) {Path QP\\$l(s)$};
\node[modimp,minimum width=2.45cm,text width=2.35cm] (temporal) at (6.70,-1.05) {Safe windows\\Temporal homotopy};
\node[mod,minimum width=1.95cm] (speed) at (9.55,-1.05) {Speed QP\\$s(t)$};
\node[modout,minimum width=2.35cm] (out) at (12.35,-1.05) {Space--time\\trajectory};
\draw[flow] (input)--(occupancy);
\draw[flowimp] (occupancy)--(spatial);
\draw[flowimp] (spatial)--(path);
\draw[flowimp] (spatial)--(temporal);
\draw[flow] (path.south) to[out=-160,in=90] (temporal.north east);
\draw[flowimp] (temporal)--(speed);
\draw[flow] (speed)--(out);
\draw[flow,dashed] (speed.south) .. controls (9.35,-2.35) and (3.35,-2.35) ..
  node[below,font=\scriptsize]{arrival-time update} (occupancy.south);
% \node[font=\scriptsize,align=center] at (5.00,2.08) {MIKU constraint compiler};
% \node[font=\scriptsize,align=center] at (10.70,-2.78) {Existing convex solvers retained};
\end{tikzpicture}
}
\caption{MIKU constraint construction and transfer through a PVD planning stack.}
\label{fig:miku_method_flow_en}
\end{figure}

\subsection{Threat-aware margins}
\label{sec:method_threat_en}

For obstacle $i$, MIKU combines collision urgency, lateral overlap, relative motion, object type, and local interaction density into a normalized threat score:
\begin{equation}
 \Theta_i=w_1f_{\mathrm{TTC}}(i)+w_2f_{\mathrm{overlap}}(i)+w_3f_{\mathrm{vel}}(i)+w_4f_{\mathrm{type}}(i)+w_5f_{\mathrm{inter}}(i).
 \label{eq:threat_score_en}
\end{equation}
All factors lie in $[0,1]$, and the weights are $(0.30,0.20,0.15,0.10,0.25)$. When the ego vehicle is behind an obstacle and closing, $\mathrm{TTC}_i=(s_i-s_{\mathrm{ego}})/(\dot{s}_{\mathrm{ego}}-\dot{s}_i)$; otherwise it is $+\infty$. The piecewise-linear TTC factor equals one for $\mathrm{TTC}_i\leq2.0$ s, zero for $\mathrm{TTC}_i\geq7.0$ s, and decreases linearly in between \citep{minderhoud2001extended}. The overlap and relative-speed factors are
\begin{align}
 f_{\mathrm{overlap}}(i)&=\min\!\left(1,\frac{\max(0,\min(l_i^+,l_{\mathrm{ego}}^+)-\max(l_i^-,l_{\mathrm{ego}}^-))}{\max(0.1,\min(W_i,W_{\mathrm{ego}}))}\right),\\
 f_{\mathrm{vel}}(i)&=\sigma(v_{\mathrm{rel},i}/v_{\max}),\qquad \sigma(x)=\frac{1}{1+e^{-5x}}.
\end{align}
Here $W_i=l_i^+-l_i^-$ is the obstacle width; the cap keeps the overlap factor in $[0,1]$ for obstacles wider than the ego vehicle.
The type factors for vulnerable road users, motor vehicles, unknown movable objects, static obstacles, and small markers are respectively $1.0$, $0.7$, $0.5$, $0.3$, and $0.15$. Density uses a 10 m neighborhood and is zero for a single obstacle. The extra margin is
\begin{equation}
 d_{\mathrm{buf},i}=d_{\mathrm{buf,min}}+(d_{\mathrm{buf,max}}-d_{\mathrm{buf,min}})\Theta_i,
 \label{eq:dynamic_delta_en}
\end{equation}
where $[d_{\mathrm{buf,min}},d_{\mathrm{buf,max}}]=[0.10,0.40]$ m. This term changes the boundary value only; it does not change the sorting or scan structure.

\subsection{Continuous lateral partition and maximum gap}
\label{sec:method_band_en}

Obstacles whose longitudinal intervals overlap are assigned to the same connected component. Sorting by the lower station edge and maintaining the largest upper edge constructs all components in $O(n\log n)$. Within one component, sort the active intervals by $u_i$ (ties by descending $v_i$). The following lemma limits the side-assignment search.

\begin{lemma}[Continuous optimal partition]
\label{lem:continuous_partition_en}
For the center-band objective in (\ref{eq:center_band_en}), there is an optimal assignment in which the obstacles sorted by $u_i$ form a prefix assigned to $L$ and a suffix assigned to $R$. Hence only $k+1$ split points need be considered.
\end{lemma}

\emph{Proof.} If both sides are nonempty, let $\beta=\min_{i\in\mathcal{R}}u_i$. Move every obstacle currently in $\mathcal{L}$ with $u_i\geq\beta$ to $\mathcal{R}$. The upper bound remains unchanged because each moved lower edge is no smaller than $\beta$, whereas the lower bound cannot increase after removing elements from $\mathcal{L}$. The width therefore does not decrease. Repeating this operation yields a continuous split. Empty-side cases give $p=0$ or $p=k$. Descending $v_i$ resolves ties in $u_i$. \hfill$\Box$

For split $p$, define the prefix maximum and candidate upper edge by
\begin{equation}
 V_0=L^-,\qquad V_p=\max\left(L^-,\max_{1\leq q\leq p}v_{(q)}\right),
 \label{eq:prefix_v_en}
\end{equation}
\begin{equation}
 U_p=\begin{cases}\min(L^+,u_{(p+1)}),&p<k,\\L^+,&p=k,\end{cases}\qquad g_p=U_p-V_p.
 \label{eq:gap_def_en}
\end{equation}

\begin{theorem}[Maximum-gap optimality]
\label{thm:max_gap_en}
For the fixed-section interval model, the optimal center-band width is $W^*=\max_{p=0,\ldots,k}g_p$. The candidates are obtained by sorting and one prefix scan, with complexity $O(k\log k)$.
\end{theorem}

\emph{Proof.} For split $p$, (\ref{eq:center_band_en}) gives the lower edge $V_p$ and the tightest upper edge $U_p$, so its width is $g_p$. Lemma~\ref{lem:continuous_partition_en} guarantees an optimal continuous split; enumerating the $k+1$ splits therefore attains the optimum. \hfill$\Box$

The theorem is exact only for this fixed-section interval abstraction. It does not certify a continuous global optimum for the full nonconvex trajectory problem. MIKU keeps the $Q=3$ widest positive bands in each component and uses dynamic programming to balance band width against lateral-center changes. If band $B_{jq}$ has center $c_{jq}$ and width $g_{jq}$ in component $j$, the sequence minimizes
\begin{equation}
 -\sum_jg_{jq_j}+\lambda_s\left(|c_{1q_1}-l_0|+\sum_{j>1}|c_{jq_j}-c_{j-1,q_{j-1}}|\right),
 \label{eq:spatial_homotopy_en}
\end{equation}
with the recurrence
\begin{equation}
 D_{jq}=-g_{jq}+\min_r\{D_{j-1,r}+\lambda_s|c_{jq}-c_{j-1,r}|\}.
 \label{eq:spatial_dp_en}
\end{equation}

\begin{figure}
\centering
\adjustbox{max width=0.82\textwidth}{\begin{tikzpicture}[x=1cm,y=1cm]
  \fill[bglight] (0,-.45) rectangle (10,.45);
  \draw[gray,thick] (0,0)--(10,0);
  \foreach \x/\name in {1.2/$O_1$,3.0/$O_2$,6.0/$O_3$,8.1/$O_4$}{
    \fill[dangerred!70] (\x-.38,-.25) rectangle (\x+.38,.25);
    \node[font=\scriptsize] at (\x,.7) {\name};
  }
  \draw[deepblue,very thick,<->] (3.38,-.85)--(5.62,-.85);
  \node[deepblue,font=\scriptsize] at (4.5,-1.15) {$g_2=U_2-V_2$ (maximum gap)};
  \foreach \x/\lab in {0/$p=0$,2.62/$p=1$,4.62/$p=2$,6.62/$p=3$,8.48/$p=4$}{
    \draw[gray!70,dashed] (\x,-.48)--(\x,.48);
    \node[font=\tiny] at (\x,-.63) {\lab};
  }
  \node[font=\scriptsize,align=center] at (5,1.15) {$k=4$ intervals yield $k+1=5$ continuous split candidates};
\end{tikzpicture}
}
\caption{Continuous split candidates for the fixed-section maximum-gap problem. Four sorted forbidden intervals induce five candidate splits; the prefix maximum is required when interval upper edges are nested.}
\label{fig:max_gap_en}
\end{figure}

\subsection{Arrival-time-dependent projection}
\label{sec:method_corridor_en}

At station $s$, the path stage queries obstacle occupancy at the estimated ego arrival time. With the current state and a constant-acceleration approximation,
\begin{equation}
 \tau(s)=\frac{-v_{\mathrm{ego}}+\sqrt{v_{\mathrm{ego}}^2+2a_{\mathrm{ego}}(s-s_{\mathrm{ego}})}}{a_{\mathrm{ego}}},
 \label{eq:arrival_time_en}
\end{equation}
with the constant-speed limit used when $|a_{\mathrm{ego}}|<10^{-3}$. A negative discriminant is marked outside the planning horizon. In rolling replanning, $\tau(s)$ is interpolated from the previous speed solution and updated at each cycle.

The numerical scenarios use $s_i(t)=s_{i,0}+v_{s,i}t$ and $l_i(t)=l_{i,0}+v_{l,i}t$. For bounded prediction errors, the longitudinal and lateral expansion at time $t$ are
\begin{equation}
 \epsilon_{s,i}(t)=\epsilon_{s0,i}+t\epsilon_{vs,i},\qquad
 \epsilon_{l,i}(t)=\epsilon_{l0,i}+t\epsilon_{vl,i}.
 \label{eq:prediction_tube_en}
\end{equation}
The physical obstacle and ego footprint are expanded by these quantities before the ST occupancy map is formed. Nominal $\tau(s)$ selects and ranks alternatives; collision checking uses the expanded occupancy.

\subsection{Temporal homotopy and ST transfer}
\label{sec:method_inject_en}

For conflict region $c$, let $I_{cr}=[t_{cr}^{in},t_{cr}^{out}]$ be each overlap interval between the path and the predicted occupancy tube. After temporal safety expansion $\eta_c$, the safe windows are
\begin{equation}
 \mathcal{W}_c=[0,T]\setminus\bigcup_r[t_{cr}^{in}-\eta_c,t_{cr}^{out}+\eta_c].
 \label{eq:safe_windows_en}
\end{equation}
Intersecting these windows with the earliest reachable arrival time yields $W_{cq}=[\alpha_{cq},\beta_{cq}]$. For conflict points ordered by station, connect adjacent nodes when
\begin{equation}
 \hat t_c\geq\hat t_{c-1}+\frac{s_c-s_{c-1}}{v_{\max}},\qquad \hat t_c\in W_{cq}.
 \label{eq:temporal_edge_en}
\end{equation}
Node cost measures deviation from $\tau(s_c)$, and edge cost discourages unnecessary switching between before- and after-obstacle passage. A bounded-width dynamic program selects one causal sequence.

For a selected window, its semantics enter the speed QP as linear box constraints:
\begin{equation}
 \begin{aligned}
 s_j^{ub}&=\min(s_j^{ub,0},s_c^- -\epsilon_s),&&t_j<\alpha_{cq},\\
 s_j^{lb}&=\max(s_j^{lb,0},s_c^+ +\epsilon_s),&&t_j\geq\beta_{cq}.
 \end{aligned}
 \label{eq:window_to_st_en}
\end{equation}
If no complete window sequence is reachable, the planner imposes a stopping upper bound before the first unreachable conflict. The downstream speed QP remains a convex QP with the same objective and kinematic equalities.

\subsection{Algorithm and complexity}

\begin{algorithm}
\small
\caption{MIKU constraint compilation}
\label{alg:miku_en}
\KwIn{Road bounds, ego state, and predicted trajectories for $n$ obstacles}
\KwOut{Center path bounds and ST position bounds}
Estimate $\tau(s)$ and compute threat-dependent center intervals for each obstacle.\;
Sort station intervals and form longitudinal connected components.\;
\ForEach{component}{
  Sort by $(u_i,-v_i)$, compute prefix maxima, and retain the $Q$ widest positive bands.\;
}
Select a continuous spatial-homotopy sequence by (\ref{eq:spatial_homotopy_en}).\;
Write the selected side assignments into the active path bounds at each station.\;
Construct uncertainty-expanded occupancy tubes and safe windows.\;
Select a causal temporal-homotopy sequence by (\ref{eq:temporal_edge_en}) and compile it into (\ref{eq:window_to_st_en}).\;
If enabled, update $\tau(s)$ from the speed solution and repeat until the iteration limit or convergence.\;
\Return{Path and ST bounds}
\end{algorithm}

For $K$ path samples, $m$ components, at most $Q$ spatial candidates per component, and a bounded temporal beam width $B$, sorting and grouping cost $O(n\log n)$. With $K_s$ retained prefixes per terminal spatial band, the layered spatial dynamic program costs $O(mK_sQ^2\log(K_sQ))$; temporal beam expansion and ranking cost $O(hBR\log(BR))$ for $h$ conflict points and $R$ reachable windows per layer. The implementation uses $Q=3$, $K_s=3$, and $B=8$. Occupancy queries are $O(nK)$ and naive interaction-density computation is $O(n^2)$ in the worst case.

\section{Experimental design}
\label{sec:experimental_design}

The experimental evaluation investigates planning performance across three distinct tiers: large-scale open-loop statistical simulation, closed-loop rolling replanning, and production software integration within Apollo Planning. Statistical experiments use deterministic scenario generators and fixed random seeds so that every method receives identical inputs, with collision and clearance evaluated against separately maintained ground-truth trajectories.

\subsection{Protocols and scenarios}
\label{sec:protocols_scenarios}

Each scenario uses a fixed horizon $T\in\{9,10,12,14\}$~s assigned by family: 9~s for crossing pedestrians and delayed crossings; 10~s for vehicle cut-ins, parked-with-oncoming traffic, and prediction noise; 12~s for narrow multi-obstacle situations; and 14~s for interleaved dynamic traffic. An episode times out when the clock reaches $T$ without achieving $s_{\mathrm{end}}\geq s_{\max}-1$~m. Four disjoint protocols are used.

\emph{Open-loop evaluation over 3,500 scenarios.} 500 seeds in each of seven families: crossing pedestrians, vehicle cut-ins, parked with oncoming traffic, narrow multi-obstacle, interleaved dynamic, prediction noise, delayed crossings. Parameters are sampled within family ranges. Planning and truth scenarios coincide except in prediction-noise, where bounded errors are injected only into planning input.

\emph{Fixed-section geometry over 4,000 instances.} 20 deterministic seeds with 200 fixed-section interval instances per seed; each has zero to eight possibly overlapping, nested, or road-crossing forbidden intervals. The scan is compared with exhaustive enumeration of all $2^k$ side assignments on center-band width.

\emph{Rolling-replanning evaluation over 700 scenarios.} 100 seeds per family. The ego vehicle is initialized from truth state $(s_0,l_0,v_0,a_0)$ at $t=0$ with a cold-start first cycle that does not use warm-started QP duals or committed homotopy. Every 0.5~s the first planned segment executes, obstacle states update, and the planner is called on shifted horizon $\max(2.0\text{ s},T-t)$ until the vehicle achieves arrival at $s\geq s_{\max}-1$~m or reaches $T$.

\emph{Finite-grid joint reference over 70 scenarios.} Stratified with 10 seeds per family and grid spacings $\Delta t=0.5$ s, $\Delta s=0.5$ m, $\Delta l=0.25$ m, $\Delta v=1.0$ m/s, beam width 20,000.

\subsection{Baselines and implementation}
\label{sec:baselines_impl}

The matched PVD baselines are defined by information exchanged at the constraint interface:
\begin{itemize}
  \item B0 is time-blind and uses the fixed local side policy.
  \item B1 adds arrival-time-dependent projection but retains obstacle-wise greedy side selection.
  \item B2 allows three arrival-time refinement iterations in the B1 policy.
  \item MIKU adds group-level continuous partition, Top-$K$ spatial homotopy, threat-aware margins, uncertainty-expanded occupancy, temporal homotopy, and at most two planning iterations with alternating updates.
\end{itemize}
All four share identical path and speed QP objectives, kinematic constraints, input, and metric implementation. Both piecewise-jerk QPs use OSQP at default tolerances with no per-call wall-clock cap on the 3,500-scenario runs. Baseline B2 uses the same QPs inside three arrival-time refinement iterations. Baseline B3 is a deterministic beam-search dynamic program on the $(t,s,l,v)$ grid with a beam width of 20,000. Its per-transition cost weights speed deviation from $v_0$, longitudinal acceleration, lateral offset, and lateral acceleration minus a progress reward; transitions that are kinematically inadmissible or collide with uncertainty-expanded obstacles are pruned before ranking. MIKU hyperparameters were fixed before evaluation and are reported completely in Table~\ref{tab:fixed_parameters}. Runs used an Intel i9-14900HX with 64 GB RAM. Apollo integration confirms compatibility with Apollo 11.0 Planning.

\begin{table}[t]
\centering
\caption{Fixed planning and uncertainty parameters used in all reported experiments.}
\label{tab:fixed_parameters}
\small
\begin{tabular}{lll}
\toprule
Parameter & Value & Role \\
\midrule
$(w_1,\ldots,w_5)$ & $(0.30,0.20,0.15,0.10,0.25)$ & threat-score weights \\
TTC breakpoints & $2.0,7.0$ s & urgency normalization \\
$(\delta_{\min},\delta_{\max})$ & $(0.10,0.40)$ m & adaptive obstacle buffer \\
$Q,K_s,\lambda_s$ & $3,3,0.05$ & spatial bands, Top-$K$, transition weight \\
$B$ & $8$ & temporal-graph beam width \\
$\eta_c$ & $0.10+0.05\min(|v_l|,2)$ s & occupancy-window expansion \\
$\epsilon_s$ & $0.05$ m & longitudinal ST guard \\
$(\epsilon_{s0},\epsilon_{l0})$ & $(0.35,0.18)$ m & initial prediction-error bounds \\
$(\epsilon_{vs},\epsilon_{vl})$ & $(0.30,0.12)$ m/s & prediction-error growth bounds \\
Path discretization & $0.5$ m & path-bound sampling \\
Speed discretization & $0.1$ s & speed-QP sampling \\
Planning iterations & at most $2$ & initial solve plus refinement \\
\bottomrule
\end{tabular}
\end{table}

\subsection{Outcomes and statistical analysis}
\label{sec:outcomes_stats}

The primary outcome is collision-free goal reaching: a run succeeds when collision-free over the horizon and reaching $s_{\mathrm{end}}\geq s_{\max}-1$ m. Collision rate uses truth rectangles; degraded-stop rate records runs that stop or fail to reach the goal without collision. Progress ratio, average speed, jerk RMS, minimum signed clearance, minimum longitudinal TTC, penalised travel time, and planning latency give safety, comfort, and computational context. Penalised travel time equals arrival time for collision-free arrivals and family-specific $T$ from Section~\ref{sec:protocols_scenarios} plus 5~s otherwise. Runtime includes the full method call; B2 includes all coordination iterations.

For each fixed seed we compute MIKU minus baseline differences before aggregation. Continuous metrics use 5,000-resample paired percentile bootstrap intervals; binary outcomes use exact McNemar tests; effect sizes are paired Cohen's $d_z$. Runtime is reported at the 50th, 95th, and 99th percentiles.

\section{Results}
\label{sec:results}

\subsection{RQ1: Main comparison}
\label{subsec:results_main}

The fixed-section test found exact agreement between the $O(k\log k)$ scan and exhaustive enumeration for all 4,000 generated instances. The maximum absolute center-band-width difference was below the $10^{-12}$ numerical test tolerance. This result establishes the theorem's fixed-section claim only; it does not imply global optimality of the complete trajectory planner.

Table~\ref{tab:main_results} reports the 3,500 paired scenarios. MIKU increased collision-free goal reaching from 62.69\% for B0 to 77.46\%, a paired difference of 14.77 percentage points (95\% CI 13.54--15.97; exact McNemar $p=1.94\times10^{-133}$). Collision decreased from 2.17\% to 1.14\% (difference $-1.03$ percentage points, 95\% CI $-1.37$ to $-0.71$; $p=2.91\times10^{-11}$). Mean progress increased by 8.93 percentage points, jerk RMS decreased by 0.120 m/s$^3$, and penalized travel time decreased by 0.81 s. The aggregate mean-clearance difference was $-0.010$ m with a 95\% CI of $[-0.024,0.004]$, so the success and collision gains cannot be interpreted as a uniform clearance increase.

\begin{table}
\tbl{Results on 3,500 paired scenarios. Values are percentages for rates, means for continuous outcomes, and milliseconds for latency.}
{
\begin{tabular}{lrrrr}
\toprule
Outcome & B0 & B1 & B2 & MIKU \\
\midrule
Collision-free goal reaching (\%) & 62.69 & 52.37 & 62.09 & \textbf{77.46} \\
Collision (\%) & 2.17 & 2.14 & 2.23 & \textbf{1.14} \\
Degraded stop (\%) & 35.40 & 45.77 & 35.97 & \textbf{21.51} \\
Progress ratio (\%) & 77.52 & 71.07 & 76.62 & \textbf{86.45} \\
Jerk RMS (m/s$^3$) & 1.252 & 1.315 & 1.300 & \textbf{1.132} \\
Runtime P50 (ms) & 10.11 & 11.37 & 31.20 & \textbf{9.98} \\
Runtime P95 (ms) & 56.79 & 63.24 & 169.49 & \textbf{55.93} \\
Runtime P99 (ms) & 138.88 & 169.73 & 463.05 & 221.80 \\
\bottomrule
\end{tabular}}
\label{tab:main_results}
\end{table}

The family-level effects in Figure~\ref{fig:evidence_dashboard} show where the overall difference originates. The largest success gain occurs in the narrow multi-obstacle family (65.4 percentage points), followed by delayed crossing (26.8 percentage points). Crossing-pedestrian and vehicle cut-in families show no aggregate success difference in this protocol, which is consistent with cases in which the baseline already has a feasible temporal response. Prediction noise yields a 3.2-point success increase and a 7.0-point collision reduction, but progress decreases by 1.12 points and jerk RMS increases by 0.041 m/s$^3$.

The finite-grid reference provides a computational trade-off rather than a winner-takes-all comparison. In 70 cases, B3 reached 80.0\% of goals and MIKU 77.1\% (paired difference $-2.86$ percentage points, 95\% CI $[-11.43,5.71]$, $p=.754$). Median runtime was 1,835.43 ms for B3 and 10.44 ms for MIKU. Thus MIKU approaches the finite reference's success rate at a much lower measured cost in this small, coarse grid experiment, while the two solvers optimize different representations.

\begin{figure}
\centering
\includegraphics[width=\textwidth]{evidence_dashboard.pdf}
\caption{Evidence summary from the frozen 3,500-case protocol. Panel (a) shows paired MIKU--B0 success effects by scenario family with 95\% bootstrap intervals; (b) shows complete-planning latency distributions; (c) shows full-method minus ablation success effects; and (d) decomposes MIKU outcomes into successful arrival, safe non-arrival, and collision.}
\label{fig:evidence_dashboard}
\end{figure}

\subsection{RQ2: Mechanisms and ablations}
\label{subsec:results_ablation}

The subtractive ablations isolate mechanisms on the same 3,500 scenarios. Removing Top-$K$ spatial homotopy reduced success by 10.17 percentage points and removing the threat-aware margin reduced it by 10.74 points relative to full MIKU. Removing the temporal homotopy graph reduced delayed-crossing success from 89.6\% to 66.8\%, a 22.8-point family-specific loss. Removing the robust occupancy tube increased prediction-noise collision from 3.6\% to 11.0\%. Arrival projection and longitudinal grouping produce smaller aggregate changes, indicating that their value is concentrated in particular configurations rather than uniform across all families.

\begin{table}
\tbl{Subtractive ablations on the same 3,500 paired scenarios. Each row removes one MIKU component.}
{
\begin{tabular}{l lrrr}
\toprule
Variant & Removed component & Success (\%) & Collision (\%) & Progress (\%) \\
\midrule
A1 & Arrival-time projection & 77.34 & 1.23 & 86.73 \\
A2 & Longitudinal grouping & 76.94 & 1.14 & 85.89 \\
A3 & Top-$K$ spatial homotopy & 67.29 & 1.14 & 79.53 \\
A4 & Threat-aware margin & 66.71 & 1.14 & 78.85 \\
A5 & Temporal homotopy graph & 74.20 & 1.17 & 84.29 \\
A6 & Robust occupancy tube & 76.97 & 2.20 & 86.60 \\
A7 & Alternating refinement & 76.63 & 1.03 & 85.92 \\
MIKU & None & \textbf{77.46} & \textbf{1.14} & \textbf{86.45} \\
\bottomrule
\end{tabular}}
\label{tab:ablation_results}
\end{table}

The mechanism cases provide a geometric explanation for these effects. In the narrow-passage case, the obstacle-wise policy raises and lowers opposing path bounds until the center interval becomes empty, whereas the group-level partition retains one continuous band. In the delayed-crossing family, the temporal graph preserves the before- and after-obstacle alternatives and enforces their causal order. These examples locate the failure mechanism; they are not substitutes for the family-level estimates.

\begin{figure}
\centering
\adjustbox{max width=\textwidth}{\begin{tikzpicture}[x=0.82cm,y=0.82cm,>=Stealth]
\begin{scope}
  \fill[bglight] (-.4,-2.25) rectangle (7.8,2.25);
  \draw[gray,thick] (-.4,2.25)--(7.8,2.25) (-.4,-2.25)--(7.8,-2.25);
  \draw[dashed,gray!60] (-.4,0)--(7.8,0);
  \foreach \x in {1.5,3.5,5.5}{
    \fill[dangerred!70] (\x,.50) rectangle (\x+.7,1.65);
    \fill[dangerred!70] (\x,-1.65) rectangle (\x+.5,-.50);
  }
  \draw[deeppink,very thick] (-.4,2.25)--(1.25,2.25)--(1.25,-1.65)--(6.35,-1.65)--(6.35,2.25)--(7.8,2.25);
  \draw[deepblue,very thick] (-.4,-2.25)--(1.25,-2.25)--(1.25,1.65)--(6.35,1.65)--(6.35,-2.25)--(7.8,-2.25);
  \node[deepblue,font=\scriptsize\bfseries] at (.55,1.10) {$\mathcal{L}$};
  \node[deeppink,font=\scriptsize\bfseries] at (.55,-1.10) {$\mathcal{R}$};
  \node[dangerred,font=\scriptsize\bfseries] at (3.8,.25) {$l^+<l^-$};
  \node[font=\small\bfseries] at (3.7,-2.75) {(a) Greedy obstacle-wise bounds};
\end{scope}
\begin{scope}[xshift=9cm]
  \fill[bglight] (-.4,-2.25) rectangle (7.8,2.25);
  \draw[gray,thick] (-.4,2.25)--(7.8,2.25) (-.4,-2.25)--(7.8,-2.25);
  \draw[dashed,gray!60] (-.4,0)--(7.8,0);
  \foreach \x in {1.5,3.5,5.5}{
    \fill[dangerred!70] (\x,.50) rectangle (\x+.7,1.65);
    \fill[dangerred!70] (\x,-1.65) rectangle (\x+.5,-.50);
  }
  \fill[lightgreen,opacity=.28] (-.4,-.50) rectangle (7.8,.50);
  \draw[deeppink,very thick] (-.4,.50)--(7.8,.50);
  \draw[deepblue,dashed,very thick] (-.4,-.50)--(7.8,-.50);
  \draw[deepblue,very thick,->] (.2,0)--(1.7,-.08)--(3.3,.06)--(4.9,-.05)--(6.5,.08)--(7.4,.02);
  \node[deeppink,font=\scriptsize\bfseries] at (.55,1.10) {$\mathcal{R}$};
  \node[deepblue,font=\scriptsize\bfseries] at (.55,-1.10) {$\mathcal{L}$};
  \node[lightgreen!60!black,font=\scriptsize\bfseries] at (3.8,1.95) {center corridor};
  \node[font=\small\bfseries] at (3.7,-2.75) {(b) Group-level MIKU bounds};
\end{scope}
\end{tikzpicture}
}
\caption{Illustration of the narrow multi-obstacle mechanism. Group-level partitioning retains a continuous center band after the obstacle-side decisions are coordinated.}
\label{fig:narrow_mechanism_en}
\end{figure}

\subsection{RQ3: Robustness, side effects, and latency tails}
\label{subsec:results_robustness}

The prediction-noise subgroup demonstrates both a benefit and a cost. MIKU reduces collisions by 7.0 percentage points relative to B0, but its mean progress is 59.73\% versus 60.85\% and its jerk RMS is 2.274 versus 2.234 m/s$^3$. This is consistent with uncertainty-expanded occupancy being conservative in cases where the error bound is large relative to the available passage.

The complete-call runtime median and P95 are similar to or lower than B0, but the P99 tail is not. MIKU's P99 is 221.80 ms, compared with 138.88 ms for B0; the empirical latency curve in Figure~\ref{fig:evidence_dashboard} shows the divergence in the upper tail. B2 is slower across all reported quantiles because it always performs three coordination iterations. These results support a typical-cycle efficiency statement, not a hard real-time guarantee.

The independent rolling-replanning protocol gives a second view of accumulated behavior. Across 700 cases, success increases from 62.29\% to 72.29\% (paired difference 10.0 percentage points, 95\% CI 7.57--12.57), while collision remains 1.29\% for both methods. Mean progress increases from 82.55\% to 90.25\%, but episode-runtime P95 increases from 827.20 to 1,361.51 ms. Because this protocol executes and replans rather than evaluating one open-loop call, it is reported as robustness evidence and not merged with Table~\ref{tab:main_results}.

\subsection{RQ4: PVD software-interface evidence}
\label{subsec:results_apollo}

The Apollo 11.0 implementation passes the generated path and ST constraints through the existing path-boundary and ST-boundary interfaces while retaining the downstream QP objectives and kinematic equalities. Three representative Dreamview cases show dynamic overtaking, narrow-passage traversal, and single-lane diversion. These cases demonstrate that the abstract constraint compilation has a compatible software insertion point and produces the intended qualitative behavior. They do not establish physical-road safety, native Apollo latency, hardware-in-the-loop performance, or complete end-to-end correctness.

\begin{figure}
\centering
\begin{minipage}{0.32\textwidth}\centering\includegraphics[width=\linewidth]{dreamview/scn01_during.png}\par\small Dynamic overtaking\end{minipage}\hfill
\begin{minipage}{0.32\textwidth}\centering\includegraphics[width=\linewidth]{dreamview/scn02_during.png}\par\small Narrow passage\end{minipage}\hfill
\begin{minipage}{0.32\textwidth}\centering\includegraphics[width=\linewidth]{dreamview/scn03_during.png}\par\small Lane diversion\end{minipage}
\caption{Representative Apollo Dreamview interface cases. The screenshots are qualitative integration evidence and are not used to compute the statistical results above.}
\label{fig:apollo_cases}
\end{figure}

\section{Discussion}
\label{sec:discussion}

\subsection{Physical interpretation of performance gains}

The 3,500-case paired comparison directly connects theoretical failure mechanisms to concrete quantitative improvements. MIKU raises collision-free goal reaching from 62.69\% to 77.46\% against baseline B0, yielding a paired increase of 14.77~pp with a 95\% confidence interval of 13.54 to 15.97. This performance gain is concentrated in configurations governed by severe spatial or temporal bottlenecks:
under narrow multi-obstacle geometry, where greedy obstacle-wise decisions cross each other and collapse opposing path bounds such that $l^+ < l^-$, MIKU's group-level partition restores a continuous, feasible center corridor, increasing success by 65.4 pp;
under delayed-crossing geometry, where time-blind bounds falsely treat future arrival-time occupancies as immediate blockages, arrival projection and temporal homotopy preserve reachable passing windows before and after the crossing obstacle, increasing success by 26.8 pp.
In contrast, in scenarios naturally characterized by a single dominant corridor, such as simple pedestrian crossings and high-speed vehicle cut-ins, performance remains consistent with the baseline, verifying that MIKU specifically targets and resolves diagnosed interface bottlenecks.

Ablation experiments substantiate this mechanism attribution. Removing Top-$K$ spatial homotopy under ablation~A3 or the threat-aware margin under ablation~A4 reduces success by 10.17~pp and 10.74~pp, respectively, demonstrating that group-level lateral coordination provides the foundational clearance necessary for spatial feasibility. Removing the temporal homotopy graph under ablation~A5 selectively cuts delayed-crossing success from 89.6\% to 66.8\%, a loss of 22.8~pp that confirms its decisive role when multiple temporal passages coexist. Finally, removing the robust occupancy tube under ablation~A6 causes prediction-noise collision rates to jump from 3.6\% to 11.0\%, affirming that uncertainty-aware buffer expansion is essential for absorbing real-world tracking and perception noise \citep{yoon2024stcorridor,yangbo2023ped,pathspeedstructured2024}.

These findings align with prior decomposition-based planning literature while clarifying the locus of improvement. \citet{fan2018baidu} and \citet{zhang2020pathqp} demonstrated that piecewise-jerk QP formulations achieve production-grade smoothness and latency. Corridor-based approaches \citep{liu2017sfc,yoon2024stcorridor} showed that explicit feasible-region construction improves trajectory quality in structured environments; MIKU extends this principle by coupling spatial and temporal corridor selection through a shared homotopy graph. The iterative anchoring strategy of \citet{chen2019em} feeds timing information back between stages, yet our B2 results, reaching 62.09\% success and remaining comparable to the time-blind B0, confirm that iterative arrival-time updates alone cannot recover lateral classes discarded by an irreversible initial side decision. For traffic system designers, the practical consequence is that upgrading the constraint interface layer---without replacing the downstream QP solvers---can recover a substantial fraction of the performance gap between decomposed and joint planners while preserving the computational budget required for real-time deployment.

\subsection{Computational efficiency and engineering trade-offs}

In operational autonomous driving systems, planning algorithms must operate within strict onboard latency budgets, typically 50 to 100~ms. Across 3,500 evaluated instances, MIKU achieves a median runtime of 9.98 ms and a P95 of 55.93 ms, demonstrating that interaction-aware constraint reconstruction adds negligible overhead during nominal driving conditions.

Computational cost concentrates primarily in the distribution tail: complete-call P99 latency rises to 221.80~ms compared with 138.88~ms for B0. This tail behavior reflects the physical reality of dense multi-obstacle traffic: branch expansion in the spatial partition and temporal homotopy graph only activates when multiple obstacles simultaneously compete for shared roadway space. For production deployment, timing guards should be budgeted against the P99 tail rather than median latency. Furthermore, under prediction noise, the 7.0 pp collision reduction incurs a modest 1.12 pp progress reduction and a 0.041 m/s\textsuperscript{3} jerk increase, illustrating an intentional and controllable engineering trade-off between absolute safety margins and aggressive forward progress.

\subsection{Comparison with joint search and iterative baselines}

Under the 70-case finite-grid protocol, the joint spatiotemporal reference B3 achieves 80.0\% success compared with 77.1\% for MIKU, with a paired difference of $-2.86$~pp, a 95\% confidence interval of $[-11.43,5.71]$, and $p=0.754$. B3's median runtime is 1,835.43 ms compared with MIKU's 10.44 ms, so the finite-grid search trades a modest success increment for a much higher computational cost.

Conversely, the iterative PVD baseline B2 executes three sequential path--speed iterations, requiring a median runtime of 31.20~ms, yet achieves only 62.09\% success, which is comparable to B0's 62.69\% and 15.37~pp lower than MIKU.

\subsection{Comprehensive insights from multi-protocol evaluation}

The multi-tiered experimental framework evaluates planning performance across distinct operational dimensions. The open-loop 3,500-scenario protocol provides statistical power and mechanism isolation under controlled environmental distributions. The 700-scenario closed-loop rolling replanning protocol demonstrates that open-loop advantages reliably accumulate over extended operational horizons, lifting overall goal-reaching success from 62.29\% to 72.29\%, an improvement of 10.0~pp, while maintaining trajectory stability across successive planning cycles.

Finally, integration into Apollo 11.0 validates architectural compatibility on an established production platform. Across dynamic overtaking, narrow passage navigation, and lane diversion scenarios in Dreamview, MIKU successfully drives the vehicle without modifying existing quadratic programming solvers or control interfaces.

\section{Limitations}
\label{sec:limitations}

The evidence has five explicit boundaries. First, the main scenarios are generated on a straight road with constant-velocity obstacle models. This may overstate the regularity of the predicted occupancy and does not test curvature, lane topology, weather, or incident effects; the next study should replay public and field-collected trajectories across road geometries. Second, the prediction-noise model uses bounded position and velocity errors. Its collision result therefore applies to the stated error envelope, not to an arbitrary predictor; calibrated multi-modal prediction distributions are needed next. Third, B0--B2 are internal matched PVD variants and B3 is a finite coarse-grid reference. The paper makes no public SOTA or fair cross-dataset claim; a matched implementation of a recent external planner is the next comparison. Fourth, Apollo evidence covers interface compatibility and representative Dreamview behavior, not native Apollo performance, hardware-in-the-loop execution, physical-road testing, or an end-to-end safety proof. Native closed-loop instrumentation and controlled-road tests are required before deployment claims. Fifth, the maximum-gap theorem is exact for a fixed-section interval abstraction and the retained candidate set. It does not certify a continuous global optimum of the full nonconvex trajectory problem; a future formulation could study completeness under changing obstacle order.

\section{Conclusion}
\label{sec:conclusion}

This paper presented MIKU, an interaction-aware constraint compiler for path--velocity planning with dynamic multi-obstacle traffic. Arrival-time-dependent projection, continuous group-level partitioning, uncertainty-expanded occupancy, and temporal homotopy convert lost interaction information into SL and ST bounds while preserving the downstream path and speed QPs. The fixed-section theorem reduces side-assignment search to an $O(k\log k)$ scan, which matched exhaustive enumeration on 4,000 deterministic instances. In 3,500 paired scenarios, collision-free goal reaching increased from 62.7\% to 77.5\% and collision decreased from 2.2\% to 1.1\% relative to the time-blind baseline, with the largest effects in narrow and delayed-crossing configurations. Typical latency remained near 10 ms, but the P99 tail was higher and rolling replanning increased episode-runtime P95, so the result supports bounded simulation and interface claims rather than a hard real-time or road-deployment guarantee. The next study should validate the same constraint transfer with calibrated public or field trajectories and native runtime instrumentation.

\section*{Funding}
\ifdefined\anonymousversion
This work was supported by public research programs; the funding sources and grant numbers will be restored in the accepted version. The funders had no role in study design, data collection, analysis, or the decision to publish.
\else
This work was supported by the Xiangyang Municipal Key Laboratory of Heterogeneous Big Data, the Hubei Superior and Distinctive Discipline Group of ``New Energy Vehicle and Smart Transportation,'' and the Key Research and Development Program of Hubei Province (Grant No. 2025BEB002).

\fi

\section*{Disclosure of interest}
The authors report there are no competing interests to declare.

\bibliographystyle{apacite}
\bibliography{references}

\end{document}
